\documentclass[aspectratio=169,10pt,english]{beamer}

\input{../../../_preambulo_beamer.tex}
\input{../../../_comandos_beamer.tex}

\selectlanguage{english}

\title{MA1001B - Statistical Modeling for Decision Making}
\subtitle{Lesson 43: Experimental Design: Factors and Randomization}
\author{}
\institute{Tecnol\'ogico de Monterrey}
\date{}

\begin{document}

\begin{frame}[plain]
  \vspace{1.2cm}
  {\Large\bfseries MA1001B - Statistical Modeling for Decision Making\par}
  \vspace{0.7cm}
  {\large Lesson 43: Experimental Design: Factors and Randomization\par}
  \vspace{0.7cm}
  {\color{TecRojo}\rule{\textwidth}{1.2pt}}
  \vspace{0.8cm}

  MA1001B Faculty\\[0.3cm]
  Term\\[0.3cm]
  Tecnol\'ogico de Monterrey
\end{frame}

\begin{frame}{The Design Question}
  \begin{alertblock}{Guiding question}
    When three packing methods are compared, what makes a mean difference a
    treatment effect rather than a who-chose-what difference?
  \end{alertblock}

  \vspace{0.6em}
  By the end of this lesson, you can:
  \begin{itemize}
    \item name the factor, its levels, the units, and the response;
    \item randomly assign 36 stations to three treatments;
    \item read group means after randomization;
    \item explain why self-selected groups are not a randomized experiment.
  \end{itemize}

  \vfill
  \scriptsize All stations, experience values, and cycle times used today are synthetic. No real company data are used.
\end{frame}

\begin{frame}{A Three-Level Factor and 36 Units}
  \small
  \begin{center}
    \begin{tabular}{lcc}
      \toprule
      \textbf{Design piece} & \textbf{In this experiment} & \textbf{Role} \\
      \midrule
      Factor & Packing method & What we assign \\
      Levels & Standard, Guided, Automated & $k=3$ treatments \\
      Units & 36 packing stations & What we assign to \\
      Response & Cycle time (seconds) & What we measure \\
      \bottomrule
    \end{tabular}
  \end{center}

  \vspace{0.5em}
  \begin{exampleblock}{Synthetic packing line}
    Random assignment places 12 stations in each method. Operator experience
    is a pretreatment covariate. It is not used to choose the method.
  \end{exampleblock}
\end{frame}

\begin{frame}[fragile]{Step 1: Random Assignment}
  \begin{columns}[T]
    \begin{column}{0.40\textwidth}
      \small
      \begin{lstlisting}
labels = np.repeat(
  treatments, 12
)
rng.shuffle(labels)
      \end{lstlisting}

      \begin{block}{Verified output}
        $n=12$ stations per method\\
        Mean experience $4.618$, $5.943$, $5.666$\\
        Auto-Standard gap $1.047694$ years
      \end{block}
    \end{column}
    \begin{column}{0.60\textwidth}
      \centering
      \includegraphics[width=\textwidth]{../figures/experimental_design_01_factor_randomization.png}
      \vspace{0.2em}
      \scriptsize Balanced assignment and covariate check from Step 1
    \end{column}
  \end{columns}
\end{frame}

\begin{frame}{Randomization Makes the Comparison a Treatment Contrast}
  \small
  True additive effects on cycle time:
  \[
    0\ \text{(Standard)},\quad -2\ \text{(Guided)},\quad -7\ \text{(Automated)}.
  \]
  After seed-42 random assignment, the observed means are $51.871576$,
  $50.039332$, and $43.852371$ seconds.
  \[
    \bar y_{\text{Automated}}-\bar y_{\text{Standard}}=-8.019206.
  \]
  The gap is near the true Automated effect of $-7$. Randomization, not a
  later adjustment, is what justifies reading it as a method comparison.
\end{frame}

\begin{frame}[fragile]{Step 2: Randomized Cycle Times}
  \begin{columns}[T]
    \begin{column}{0.40\textwidth}
      \small
      \begin{lstlisting}
y = (
  52 + effect
  - 0.35 * (
    exp - exp.mean()
  ) + noise
)
      \end{lstlisting}

      \begin{block}{Verified output}
        Means $51.872$, $50.039$, $43.852$\\
        Randomized gap $-8.019206$\\
        True Automated effect $-7$
      \end{block}
    \end{column}
    \begin{column}{0.60\textwidth}
      \centering
      \includegraphics[width=\textwidth]{../figures/experimental_design_02_randomized_outcomes.png}
      \vspace{0.2em}
      \scriptsize Randomized cycle times from Step 2
    \end{column}
  \end{columns}
\end{frame}

\begin{frame}{Step 3: Self-Selection Is Not a Treatment}
  \begin{columns}[T]
    \begin{column}{0.42\textwidth}
      \small
      Self-selected mean experience: $2.678$, $5.655$, $7.895$ years.

      \vspace{0.4em}
      Observational Auto-Standard $Y$ gap: $-14.809630$.

      \vspace{0.4em}
      Experience gap $5.216727$ versus randomized $1.047694$.

      \vspace{0.5em}
      \textbf{Limit:} an observational factor is not a randomized treatment.
      Lesson 44 analyzes the randomized three-group response with ANOVA.
    \end{column}
    \begin{column}{0.58\textwidth}
      \centering
      \includegraphics[width=\textwidth]{../figures/experimental_design_03_self_selection_limit.png}
      \vspace{0.2em}
      \scriptsize Randomized versus self-selected means from Step 3
    \end{column}
  \end{columns}
\end{frame}

\begin{frame}{Concept Check}
  \begin{enumerate}
    \item What is the experimental factor, and what are its levels?
    \item Why does random assignment of 36 stations matter for experience?
    \item The randomized Auto-Standard gap is $-8.019$ seconds and the true
      effect is $-7$. Why can they differ?
    \item Why is the observational gap $-14.810$ not a treatment effect?
  \end{enumerate}

  \vspace{0.6em}
  \begin{block}{Connection to Lesson 44}
    The session can continue directly into one-way ANOVA, sums of squares,
    and the $F$ ratio for the three randomized means.
  \end{block}

  \vfill
  \scriptsize Source: OpenStax, \textit{Introductory Business Statistics 2e},
  Chapter 12, Section 12.2. CC BY-NC-SA.
\end{frame}

\appendix



\end{document}
